\chapter{Evaluation of $\|\be\|_{\hds}$ for the Taylor Green Vortex}
\label{app:tg:seminorm:general}

{\color{red}
In this Appendix we {\color{red}compute} the $\hds$ seminorm for the Taylor Green Vortex. The $\hdo$ seminorm is required for properly scaling the initial conditions of the Euler and Euler-Voigt equations and the $\hdt$ is required for comparing our results to those of \cite{zhao_systematic_2023}.
\\~\\
The 3D Taylor Green Vortex on the 3D torus $\Omega$ with characteristic length $L$ is given in (\ref{eq:TGV}). To simplify the notation we drop the subscript $TGV$. The formula for $\|\be\|_{\hds}$ is

\begin{equation*}
\|\be\|_{\hds}^2 := \sum_{|\alpha|=s}\frac{s!}{\alpha!}\|d^\alpha\be\|_{L^2}^2.
\end{equation*}
This can be expressed in integral form as
\begin{equation*}
\|\be\|_{\hds}^2 = \sum_{|\alpha|=s}\frac{s!}{\alpha!}\ls \sum_{i=1}^3 \int_{\Omega} \lr \frac{\partial^{|\alpha|}}{\partial_1^{\alpha_1}\partial_2^{\alpha_2}\partial_3^{\alpha_3}} \eta_i \rr^2\,d\x \rs.
\end{equation*}
For the TGV, the third term $\eta_3$ is zero and the partial derivative of the other two terms $\eta_1$ and $\eta_2$ are
\begin{equation*}
\frac{\partial^{|\alpha|}}{\partial_1^{\alpha_1}\partial_2^{\alpha_2}\partial_3^{\alpha_3}} \eta_i = V_0\lr\frac{2\pi}{L}\rr^{|\alpha|}f(x_1)g(x_2)h(x_3), \qquad i=1,2,
\end{equation*}
where $f$, $g$, and $h$ are either $\sin(2\pi x_j/L)$ or $\cos(2\pi x_j/L)$. This allows us to simplify the equation to
\begin{equation*}
\|\be\|_{\hds}^2 = \sum_{|\alpha|=s}\frac{s!}{\alpha!}\ls \sum_{i=1}^2 V_0^2\lr\frac{2\pi}{L}\rr^{2|\alpha|} \int_{\Omega}  f(x_1)^2g(x_2)^2h(x_3)^2\,d\x \rs.
\end{equation*}
Because the integrand is separable, the integral is equal to the product of,
\begin{equation*}
\int_{0}^L \cos^2(2\pi x_j/L)\, dx_j = \frac{L}{2}, \qquad \text{and} \qquad \int_{0}^L \sin^2(2\pi x_j/L)\, dx_j = \frac{L}{2}.
\end{equation*}
Therefore, the equation can be simplified to
\begin{equation*}
\|\be\|_{\hds}^2 = \sum_{|\alpha|=s}\frac{s!}{\alpha!}\ls 2V_0^2\lr\frac{2\pi}{L}\rr^{2|\alpha|} \lr\frac{L}{2}\rr^3 \rs = \sum_{|\alpha|=s}\frac{s!}{\alpha!}\ls 2^{2|\alpha|-2}V_0^2\pi^{2|\alpha|}L^{3-2|\alpha|} \rs.
\end{equation*}
For $s=1$, there are three tuples $\alpha$ with degree $|\alpha|=1$. These are the three permutations of $\alpha=(1,0,0)$ with $\alpha! = 1$. Therefore, we obtain the result
\begin{equation*}
\|\be\|_{\hdo} = \sqrt{3L}V_0\pi.
\end{equation*} 
For $s=3$, there are ten tuples $\alpha$ with degree $|\alpha|=3$. These are $\alpha=(1,1,1)$ with $\alpha!=1$, the three permutations of $\alpha=(3,0,0)$ with $\alpha!=6$, and the six permutations of $\alpha=(1,2,0)$ with $\alpha!=2$.  Therefore,
\begin{equation*}
\|\be\|_{\hdt}^2 = \ls \frac{3!}{1} + 3\frac{3!}{6} + 6\frac{3!}{2}\rs 2^{4}V_0^2\pi^{6}L^{-3}.
\end{equation*} 
From this, we obtain the result
\begin{equation*}
\|\be\|_{\hdt} = \frac{12\sqrt{3}V_0\pi^3}{\sqrt{L^3}}.
\end{equation*} 
 
\noindent In Appendices \ref{app:time:scale:euler} and \ref{app:time:scale:euler:voigt} we derived the scaling property of the Euler and Euler-Voigt equations, respectively. This property requires that $\beta/\lambda\gamma = 1$, where $\beta$ is the scaling factor in space, $\lambda$ is the scaling factor in time, and $\gamma$ is the scaling factor of the 'size' of the solution. The optimization problem of \cite{zhao_systematic_2023} searched for optimal initial conditions with a fixed $\|\be\|_{\hdt} = 1$ on a domain with 3D periodic cube with side length $L=1$. Therefore, in this model, $V_0 = 1/(12\sqrt{3}\pi^3)$. Let $\beta=\gamma=\lambda = 1$.   
\\~\\
Our optimization problem searches for optimal initial conditions with a fixed $\|\be\|_{\hdo} = \sqrt{3}/2$. Our search is also performed on a domain with $L=1$. Using these results, we obtain the relative scaling of our model to that of \cite{zhao_systematic_2023}. These are $\beta = 1$ and $\gamma=1/6\sqrt{3}pi^2$. Thus, to satisfy the scaling property, $\gamma = 6\sqrt{3}\pi^2$. Therefore, one unit of time in our model is equivalent to $6\sqrt{3}\pi^2$ units of time in the model of \cite{zhao_systematic_2023}.


}

